\section{Current Doctrine, Canon Law, and Sacramental Discipline}

\subsection{Legal scope block}

\begin{massbox}{Law and jurisdiction applied}
\textbf{Governing law:} the 1983 \emph{Code of Canon Law} for the Latin Church, official Latin text and the Holy See's identified English translation, including Book VI as revised by Francis's apostolic constitution \emph{Pascite gregem Dei} (23 May 2021), effective 8 December 2021.\\[0.35em]
\textbf{Persons and jurisdiction:} Latin Catholics universally, subject to the Code's rules on applicability. The Code does not govern Eastern Catholic Churches automatically (CIC can.~1); an Eastern Catholic requires analysis under the CCEO, applicable particular law, and competent authority.\\[0.35em]
\textbf{Other controlling acts:} CDF declarations of 17 February 1981 and 26 November 1983; CDF doctrinal reflection of 23 February 1985; DDF response approved by Francis on 13 November 2023; for exorcism and prayer publication, CIC cann.~826, 829--830, and 1172, the CDF letter of 29 September 1985, the 2000 instruction on prayers for healing, and the current Roman Ritual.\\[0.35em]
\textbf{As-of date:} the Holy See's Code, promulgation sources, and DDF document index were checked through 13 July 2026. No later universal Roman act reversing the 1983 judgment was located.\\[0.35em]
\textbf{Assumed facts:} voluntary active enrollment in an organization that is actually Masonic. The legal result can change with the person's age, Church \emph{sui iuris}, knowledge, freedom, office, clerical or religious status, public conduct, the body's actual activity, a penal precept, particular law, and prior censures.
\end{massbox}

This section distinguishes moral doctrine, sacramental obligation, and penal law. It gives no ruling on an actual person, lodge, penalty, or admission to Communion. Those questions require complete facts and competent pastoral or canonical counsel.

\subsection{What remains universally clear}

The CDF declaration of 26 November 1983 gives four propositions:

\begin{enumerate}[leftmargin=1.6em,itemsep=0.25em]
\item The absence of the word \emph{Masonic} from the new Code reflected an editorial criterion by which associations could be contained in wider categories.
\item The Church's negative judgment remained unchanged because Masonic principles had always been considered irreconcilable with her doctrine.
\item Membership remained forbidden; faithful enrolled in Masonic associations were in a state of grave sin and could not receive Holy Communion.
\item Local ecclesiastical authorities could not issue a judgment on Masonic associations that derogated from the declaration.
\end{enumerate}

John Paul II approved the declaration and ordered its publication. The CDF's 1985 reflection explained that “state of grave sin” indicated that affiliation is objectively a grave sin. It did not make the impossible assertion that the dicastery had judged the internal knowledge and consent of every member. Catholic moral teaching requires grave matter, full knowledge, and deliberate consent for mortal sin in the complete subjective sense (CCC 1857--1861). The DDF's 2023 Philippine response accordingly spoke of those “formally and knowingly enrolled” who had embraced Masonic principles, while restating that active membership itself is forbidden and that the measures apply to clerics as well.

The distinction does not make the norm optional. Objective gravity tells conscience what may not be chosen. Reduced knowledge or freedom may lessen personal guilt; it does not turn continued informed adherence into a good. Nor may one presume reduced culpability simply because many members understand the lodge as benevolent. The pastoral task is to tell the truth in a way that permits a person to act freely, repent, and return.

\subsection{No present automatic excommunication merely for enrollment}

The 1917 Code explicitly named the Masonic sect and attached an \emph{ipso facto} excommunication reserved to the Apostolic See. The 1983 Code abrogated that Code and contains no equivalent provision attaching \emph{latae sententiae} excommunication to Masonic enrollment as such. Current canon 1314 states that a penalty is ordinarily \emph{ferendae sententiae}---binding after imposition---unless law or precept expressly makes it automatic. Current canon 1374 does not use automatic-penalty language.

It is therefore false to tell every present-day Latin Catholic Mason, solely from that membership fact, “you are automatically excommunicated under canon 1374.” The moral prohibition is not an excommunication. The inability to receive Communion stated by the CDF is not itself proof that an ecclesiastical censure has been incurred. Excommunication, interdict, objective grave sin, and a personal obligation to refrain from Communion are distinct categories even when some consequences overlap.

This correction does not weaken the Church's judgment; it states the law the Church actually promulgated. Inventing a penalty is no service to discipline. It may terrify consciences, distort sacramental reconciliation, and obscure the deeper reason for withdrawal.

\subsection{Canon 1374: association plotting against the Church}

As revised, canon 1374 reads: “A person who joins an association which plots against the Church is to be punished with a just penalty; one who promotes or takes office in such an association is to be punished with an interdict.” Three elements must not be elided.

First, the association must in fact plot against the Church. The canon uses a broader category rather than a named list. A lodge or Masonic obedience whose organized activity satisfies that element can fall under the canon. The 1983 doctrinal declaration does not mean every Masonic body is legally deemed to satisfy every factual element of this penal canon without inquiry; its point is that practical hostility is not the whole moral question.

Second, mere joining and promotion or office have different prescribed consequences. Joining attracts an indeterminate “just penalty.” Promotion or taking office attracts interdict. Because the canon does not make either penalty \emph{latae sententiae}, competent authority must proceed according to law before a penalty binds as imposed. The revised Book VI emphasizes pastoral charity, restoration of justice, reform of the offender, repair of scandal, the right of defense, and canonical equity (cann.~1311, 1341--1349).

Third, penal imputability must be established. Canon 1321 presumes innocence and requires an external violation gravely imputable by malice or culpability before punishment. Canons 1323--1325 supply excusing or mitigating circumstances, including non-culpable ignorance and reduced imputability. A general article cannot decide these facts. A bishop or judge cannot replace process with a rumor; a member cannot replace the law with the assumption that good private intention erases organized hostile action.

Other canons could become relevant on additional facts: apostasy, heresy, or schism (can.~1364); obstinate rejection of certain teaching after warning (can.~1365); incitement against ecclesiastical authority (can.~1373); prohibited participation in religious rites (can.~1381); obligations of clerics and religious; office suitability; or a particular penal precept. None should be stacked automatically onto the bare label “Mason.” Each has its own elements, competent authority, and defenses.

\subsection{Holy Communion: canons 915 and 916 are not interchangeable}

Canon 916 addresses the communicant's own obligation: one conscious of grave sin is not to celebrate Mass or receive the Lord's Body without prior sacramental confession, unless there is grave reason and no opportunity to confess, in which case perfect contrition and the resolve to confess as soon as possible are required. The CDF declaration directly enlightens conscience: enrolled faithful may not receive Communion. The 1985 reflection identifies the affiliation as objectively grave.

Canon 915 addresses the minister of Communion and those under an imposed or declared censure or “obstinately persevering in manifest grave sin.” Its application demands public and legally relevant facts: gravity, manifestation, persistence, and obstinacy. Private or uncertain membership does not automatically establish all those elements. A minister should follow competent ecclesiastical direction rather than conduct an improvised investigation at the Communion rail. Conversely, a public case with warnings and persevering active advocacy may raise external-forum questions not resolved by saying culpability is ultimately known to God.

The distinction protects both the Eucharist and persons. Canon 916 is not optional because canon 915 has not been publicly applied. Canon 915 is not mechanically triggered by every rumor of an objective grave act. Sacramental discipline is neither private self-certification nor summary public punishment.

\subsection{Confession, withdrawal, and purpose of amendment}

The sacrament of Penance reconciles those who, moved by grace, repent, confess, and intend amendment (CIC cann.~959, 987; CCC 1451--1460). A Catholic who recognizes the conflict should stop participating, seek sound pastoral counsel, and take effective steps to end membership. Depending on the body's rules and the person's situation, that may require a written resignation, return of office or property, cessation of dues and meetings, repudiation of incompatible promises, and repair of scandal or harm.

A confessor needs the relevant facts but should not demand a theatrical disclosure of every rumor or ritual detail. The moral object is continued adherence and participation. True purpose of amendment does not require certainty that withdrawal will be socially painless; it requires the sincere will to leave and prudent action toward that end. A person facing threats, employment consequences, family coercion, or danger should involve the competent pastor and, when necessary, civil counsel or safeguarding authorities. No supposed Masonic secrecy binds against reporting crime or protecting a victim.

If an old 1917-Code censure, current canon 1374 proceeding, public office, clerical status, or irregular canonical history may be involved, the confessor or pastor should consult the diocesan chancery or a qualified canonist while preserving the sacramental seal and appropriate confidentiality. The law contains faculties and procedures for remission; one should neither presume a censure nor despair of reconciliation.

\subsection{Exorcism is governed ministry, not self-help}

Canon 1172 provides that no one may lawfully pronounce exorcisms over the possessed without the local ordinary's special and express permission; that permission is given only to a priest endowed with piety, knowledge, prudence, and integrity. The reserved major rite is therefore not made available by believing that Masonry is spiritually dangerous, finding a formula online, or standing in a particular building. The Church first requires moral certainty and careful discernment, including appropriate medical, psychological, and psychiatric evaluation.

The CDF's 29 September 1985 letter applies this discipline specifically to the Leonine formula against Satan and the fallen angels. The faithful may not use that formula, much less the complete exorcism, and unauthorized assemblies may not interrogate demons or seek their identities. The letter does not forbid ordinary prayer for deliverance: it directs the faithful toward Christ's victory, the sacraments, and the intercession of Mary, the angels, and the saints. The 2000 CDF instruction further requires exorcism to remain in strict dependence upon the diocesan bishop and forbids inserting exorcistic or healing services into Mass, the sacraments, or the Liturgy of the Hours.

For Masonry this means that the objective remedy and the extraordinary remedy cannot be reversed. Ending enrollment, renouncing sinful obligations, confessing one's own grave sins, and returning to Eucharistic life address the known moral conflict. A major exorcism addresses a distinct fact---possession---that competent authority must discern. Neither a Mason nor his descendant should infer that fact from membership, ancestry, illness, resistance to a long prayer, or the alleged demon assigned to a degree.

\subsection{Publication approval and liturgical promulgation}

Canon 826 \S3 requires the local ordinary's permission to publish prayer books. Canon 829 limits approval to the original text submitted, not later editions or translations; canon 830 describes the censor's review for material contrary to faith or morals. An imprimatur is therefore meaningful ecclesiastical permission for a specified publication. It is not promulgation into the Roman Ritual, a declaration that every contingent historical or medical assertion is revealed doctrine, permission for every user to employ imperative formulas, or authentication of altered web versions.

These distinctions control references to the Ripperger collection and its derivatives. The publisher reports a Denver imprimatur for \emph{Deliverance Prayers: For Use by the Laity}, and an accessible edition identifies approval in 2018. That warrants describing the identified edition as an approved private prayer book. It does not warrant calling its Masonry prayer “the Church's official exorcism,” extending its approval to every long or short online version, or ignoring the compiler's own account that the prayer was adapted from a Protestant predecessor.

\subsection{Local authority, catechesis, and particular cases}

The 1983 reservation prevents local authorities from declaring a kind of Masonry compatible in derogation from the universal judgment. It does not disable a bishop from investigating facts, applying canon 1374, issuing legitimate particular norms within competence, directing clergy, protecting ecclesial office, or teaching the reasons for the prohibition. The 2023 DDF response expressly urged coordinated episcopal catechesis and consideration of a public statement.

Nor does the reservation prevent accurate distinctions. A bishop may need to determine whether a body is actually Masonic rather than merely fraternal, whether a named association plots against the Church for penal purposes, whether a person is enrolled, whether scandal is public, and what pastoral path is just. The universal answer to “may a Catholic actively belong to Freemasonry?” is no. The factual and canonical consequences of a concrete case still require competent judgment.

\begin{fourstable}{Question}{Universal answer}{Fact-sensitive issue}{Who should act}
May a Catholic actively enroll? & No; principles are irreconcilable and membership is forbidden. & Whether the body really is Masonic and whether enrollment is active. & The Catholic forms conscience; pastor assists; authority judges disputed facts.\\
Is the matter grave? & Yes, objectively, under the CDF's authoritative explanation. & Full knowledge, deliberate consent, duress, and other factors affect personal culpability. & Penitent and confessor in conscience; God judges the person.\\
Is there automatic excommunication now? & Not merely from enrollment under the current Latin Code. & Another delict, old-law history, or a specific precept may alter a case. & Qualified canonist and competent authority.\\
May the person receive Communion? & The declaration says enrolled faithful may not; can.~916 governs conscience. & Can.~915 requires manifest, grave, obstinate perseverance for ministerial denial outside a censure. & Communicant, confessor, pastor, and competent authority in their proper fora.\\
\end{fourstable}
